\subsection{Properties of the integral of a scalarly essentially integrable function}

PROPOSITION 5. --- Let \(\mu\) be a bounded positive measure on T, S a \(\mu\) -measurable set carrying \(\mu\) (Ch. V, §5, No.~7, Def. 4), f a scalarly \(\mu\) -integrable (*) function with values in F. Let D be the closed convex envelope of \(\mathbf{f}(\mathrm{S})\) in the space \(\mathrm{F}'^*\) equipped with the topology \(\sigma(\mathrm{F}'^*, \mathrm{F}')\) . Then \(\int \mathbf{f} d\mu \in \mu(\mathrm{T})\mathrm{D}\) .

Since D is the intersection of the closed half-spaces containing \(\mathbf{f}(\mathrm{S})\) (TVS, II, §5, No.~3, Cor. 1 of Prop. 4), it suffices to prove that the relation \(\langle\mathbf{f}(t),\mathbf{z}^{\prime}\rangle\leqslant a\) for all \(t\in S\) (where \(z^{\prime}\in F^{\prime},a\in R\) ) implies \(\langle z^{\prime},\int f d\mu\rangle\leqslant a\cdot\mu(T)\) ; but since \(\int f d\mu=\int_{S}f d\mu\) , this follows from Ch. IV, §4, No.~2, Cor. 1 of Th. 1.

COROLLARY. --- Let \(\mu\) be a bounded positive measure on T, S a \(\mu\) -measurable set carrying \(\mu\) , and f a mapping of T into F, scalarly \(\mu\) -measurable and such that \(\mathbf{f}(\mathbf{S})\) is contained in a weakly compact convex subset A of F. Then f is scalarly \(\mu\) -integrable, and \(\int f d\mu \in \mu(T)A \subset F\) .

For every \(z^{\prime} \in F^{\prime}\) , \(\langle z^{\prime}, f \rangle\) is \(\mu\) -measurable and bounded in S, hence integrable, which proves that f is scalarly integrable. Moreover, since A is compact in \(F_{\sigma}\) , it is closed in \(F^{\prime*}\) , and the closed convex envelope of \(\mathbf{f}(S)\) in \(F^{\prime*}\) is contained in A, whence the corollary.

PROPOSITION 6. --- Let f be a scalarly essentially μ-integrable function with values in F, such that \(\int f d\mu \in F\) . For every lower semi-continuous semi-norm q on F, 5

\[
q \left(\int \mathbf {f} d \mu\right) \leqslant \int^ {\bullet} (q \circ \mathbf {f}) d \mu .
\]

Let D be the set of \(z \in F\) such that \(q(\mathbf{z}) \leqslant 1\) ; D is closed, convex, and contains 0, therefore \(D = D^{\circ\circ}\) (TVS, II, §6, No.~3, Cor. 3 of Th. 1). It therefore suffices to prove that for every \(z' \in D^{\circ}\) , one has \(\left|\left\langle z', \int f d\mu \right\rangle\right| \leqslant \int^{\bullet}(q \circ f) d\mu\) ; but this follows at once from the fact that \(|\langle z', f(t) \rangle| \leqslant q(f(t))\) for every \(t \in T\) .

Note that the numerical function \(q \circ f\) need not be \(\mu\) -measurable (Exer. 12).

PROPOSITION 7. --- Let f be a mapping of T into F, scalarly essentially \(\mu\) -integrable, such that for every compact subset K of T, \(\mathbf{f}(\mathbf{K})\) is contained in a weakly compact, balanced, convex subset of F. Then \(\int f d\mu\) belongs to the bidual \(F''\) of F.

For every compact subset \(\mathrm{K}\) of \(\mathrm{T}\) ,

\[
\int \mathbf {f} \varphi_ {\mathrm{K}} d \mu = \int (\mathbf {f} \varphi_ {\mathrm{K}}) d (\varphi_ {\mathrm{K}} \cdot \mu);
\]

the Cor. of Prop. 5 can be applied to the bounded measure \(\varphi_{\mathrm{K}} \cdot \mu\) and the function \(\mathbf{f}\varphi_{\mathrm{K}}\) , consequently \(\int \mathbf{f}\varphi_{\mathrm{K}} d\mu \in \mathrm{F}\) . For every \(\mathbf{z}' \in \mathrm{F}'\) , \(\langle \mathbf{z}', \mathbf{f} \rangle\) is essentially \(\mu\) -integrable, consequently (Ch. V, §1, No.~3, Prop. 10)

\[
\int \langle {\bf z} ^ {\prime}, {\bf f} \rangle d \mu = \lim _ {\mathrm{K}} \int \langle {\bf z} ^ {\prime}, {\bf f} \rangle \varphi_ {\mathrm{K}} d \mu ,
\]

the limit being taken with respect to the increasing directed set of compact subsets of T. One concludes that, with respect to this set, \(\int f\varphi_{K}d\mu\) converges to \(\int f d\mu\) for the topology \(\sigma(F'^{*},F')\) . Now,

\[
\left| \left\langle \mathbf {z} ^ {\prime}, \int \mathbf {f} \varphi_ {\mathrm{K}} d \mu \right\rangle \right| = \left| \int \langle \mathbf {z} ^ {\prime}, \mathbf {f} \rangle \varphi_ {\mathrm{K}} d \mu \right| \leqslant \int | \langle \mathbf {z} ^ {\prime}, \mathbf {f} \rangle | d \mu ,
\]

which proves that the set of elements \(\int f\varphi_{K}d\mu\) is a bounded subset of \(F_{\sigma}\) , hence also of F (TVS, IV, §1, No.~1, Prop. 1). Proposition 7 is therefore a consequence of the following lemma:

Lemma 1. --- The closure in \(\mathbf{F}^{\prime*}\) (for the topology \(\sigma(\mathbf{F}^{\prime*},\mathbf{F}^{\prime})\) ) of every bounded subset of \(\mathbf{F}\) is contained in the bidual \(\mathbf{F}''\) .

For, a bounded subset of F is contained in the polar (in \(F''\) ) of a neighborhood of 0 in the strong dual \(F'\) of F, hence is relatively compact in \(F''\) for \(\sigma(\mathrm{F}'',\mathrm{F}')\) (TVS, III, §3, No.~5, Prop. 7 and No.~4, Cor. 2 of Prop. 4); since \(\sigma(\mathrm{F}'',\mathrm{F}')\) is induced by \(\sigma(\mathrm{F}'*,\mathrm{F}')\) , the lemma is proved.

COROLLARY. --- Suppose F is semi-reflexive, and let f be a scalarly essentially \(\mu\) -integrable mapping of T into F such that, for every compact subset K of T, \(\mathbf{f}(\mathbf{K})\) is bounded. Then \(\int f d\mu\) belongs to F.

For, every bounded subset of F is relatively weakly compact (TVS, IV, §2, No.~2, Th. 1), and \(F = F''\) .

PROPOSITION 8. --- Let \(\mu\) be a bounded positive measure on T, S a \(\mu\) -measurable set carrying \(\mu\) , f a \(\mu\) -measurable mapping of T into F, such that f(S) is contained in a complete, bounded, balanced convex subset B of F. Then, f is scalarly \(\mu\) -integrable and \(\int f d\mu \in \mu(T)B \subset F\) .

Since S is \(\mu\) -integrable, there exists a partition of S formed by a \(\mu\) -negligible set N and a sequence \((\mathrm{K}_{n})\) of compact subsets such that the restriction of f to each \(K_{n}\) is continuous (Ch. IV, §4, No.~6, Cor. 3 of Th. 4 and §5, No.~1, Def. 1); \(\mathbf{f}(\mathrm{K}_n)\) is therefore a compact subset of F. The closed, balanced convex envelope \(\mathbf{B}_n\) of \(\mathbf{f}(\mathrm{K}_n)\) is then pre-compact (TVS, II, §4, No.~1, Prop. 3) and is contained in the complete subset B of F, therefore it is compact, and a fortiori weakly compact. Consequently (Cor. of Prop. 5) \(\mathbf{f}\varphi_{\mathrm{K}_n}\) is scalarly \(\mu\) -integrable, and

\[
\mathbf {z} _ {n} = \int \mathbf {f} \varphi_ {\mathrm{K} _ {n}} d \mu \in \mu (\mathrm{K} _ {n}) \mathrm{B} _ {n} \subset \mu (\mathrm{K} _ {n}) \mathrm{B}.
\]

For every continuous semi-norm p on F, it follows that

\[
p (\mathbf {z} _ {n}) \leqslant \mu (\mathrm{K} _ {n}) \cdot \sup _ {\mathbf {x} \in B} p (\mathbf {x});
\]

since B is bounded and since the series with general term \(\mu(\mathrm{K}_{n})\) is convergent and has sum \(\mu(\mathrm{T})\) , one sees that the sequence with general term \(s_{n}=z_{1}+z_{2}+\cdots+z_{n}\) is a Cauchy sequence in the complete subset \(\mu(\mathrm{T})\mathrm{B}\) of F. This sequence therefore converges to an element s of \(\mu(\mathrm{T})\mathrm{B}\) ; since one can suppose that \(\mathbf{f}(t)=0\) on T-S, Lebesgue's theorem applied to each of the functions \(\langle z', f\rangle\) ( \(z'\in F'\) ) proves that \(s=\int f d\mu\) .
% semantic-fragments: legacy-input-seam
\relax{}
